% ===========================
% IMPLEMENTATION PLAN — BC / Top-K / LapPE / (m)HC-GNN
% ===========================
\documentclass[11pt]{article}
\usepackage[a4paper,margin=2cm]{geometry}
\usepackage{amsmath,amssymb,mathtools}
\usepackage{hyperref}
\usepackage{enumitem}
\usepackage{booktabs}

\title{Implementation Plan — Betweenness Centrality Prediction\\
Inductive Betweenness Centrality Prediction on Synthetic Graph Families}
\author{}
\date{}

\begin{document}
\maketitle
\tableofcontents

% ==========================================================
\section{Research question and scope}
% ==========================================================
\subsection{Central question}
\paragraph{RQ.}
\emph{Can deep GNNs predict betweenness centrality in an inductive setting (train on small graphs, test on larger graphs), and do multi-stream residual mechanisms (HC/mHC/mHC-lite) improve depth scaling and Top-$K$ retrieval compared to strong deep-GNN baselines?}

\subsection{What “success” means}
We consider BC prediction successful if the model:
(i) ranks nodes correctly (global correlation), and
(ii) reliably retrieves the highest-BC nodes (Top-$K$).

\subsection{Inductive-only experimental protocol}
We train on many small graphs and test on larger graphs drawn from multiple canonical families (e.g., ER/BA/SBM/WS/RGG), to avoid learning a single “graph style”.
We explicitly report in-distribution (ID) vs out-of-distribution (OOD) generalization:
ID uses the same families for train/test, while OOD uses disjoint families (e.g., train ER+BA+SBM, test WS+RGG).

\paragraph{Why HC, and why constraints.}
HC is the baseline multi-stream routing mechanism that may help depth scaling on its own.
mHC/mHC-lite are constrained variants: they keep the same routing but constrain $H^{\mathrm{res}}$ to test whether stability, not just routing capacity, is the key factor.

% ==========================================================
\section{Main hypotheses}
% ==========================================================
\paragraph{H1 (depth bottleneck).}
Increasing depth in standard MPNNs eventually degrades performance due to over-smoothing/over-squashing and optimization issues.

\paragraph{H2 (multi-stream routing helps depth).}
HC-style multi-stream residual routing improves optimization stability and preserves useful signal across depth, enabling deeper networks to outperform strong deep-GNN baselines on ranking metrics.

\paragraph{H3 (structure helps).}
Injecting structural information (e.g., LapPE, degree features) improves BC ranking and may reduce the required depth for comparable performance.

\paragraph{H4 (constraint matters at extreme depth).}
At sufficient depth, enforcing (approx./exact) doubly-stochastic constraints on $H^{\mathrm{res}}$ (mHC/mHC-lite) improves stability and ranking performance compared to unconstrained HC, while differences may be smaller at shallow depth.

% ==========================================================
\section{Task definition: labels, objectives, metrics}
% ==========================================================
\subsection{Datasets (inductive protocol)}
Train set: $M_{\\mathrm{train}}$ graphs of size $N_{\\mathrm{train}}$ (e.g., 100 graphs of 100 nodes).
Test set: graphs of size $N_{\\mathrm{test}}$ (e.g., 1000 nodes), ideally multiple graphs to estimate variance.
Families: at least five standard families from the literature: ER (Erd\H{o}s--R\'enyi), BA (Barab\'asi--Albert), SBM, WS, and RGG (or configuration model).
We generate multiple families explicitly to prevent the model from learning a single graph “style”.

\subsection{Betweenness centrality labels}
For small graphs, BC is computed exactly (Brandes algorithm).
For larger graphs, labels may be computed either exactly (if feasible, e.g., 1000 nodes depending on density) or approximately (sampling-based BC), as an engineering trade-off.

\subsection{Target transformation}
BC is heavy-tailed. We use:
\[
y(v)=\log(1+\mathrm{BC}(v)),\qquad 
\tilde y(v)=\frac{y(v)-\mu_{\mathrm{train}}}{\sigma_{\mathrm{train}}}.
\]

\subsection{Training objectives}
We optimize a ranking objective only, since Top-$K$ retrieval is the primary goal.

\paragraph{Ranking objective.}
Pairwise RankNet-style loss over sampled node pairs $(u,v)$:
\[
\mathcal{L}_{\mathrm{rank}}
=\sum_{(u,v)}\log\!\left(1+\exp\!\big(-s_{uv}(\hat y(u)-\hat y(v))\big)\right),
\quad s_{uv}=\mathrm{sign}(y(u)-y(v)).
\]

\subsection{Evaluation metrics}
We report:

\paragraph{Global ranking agreement.}
Spearman $\rho$ (and optionally Kendall $\tau$) between ranks($\hat y$) and ranks($y$).

\paragraph{Top-$K$ retrieval.}
Precision@K and NDCG@K for $K\in\{10,20,50\}$ and proportional $K$ (top 1\%, top 5\%).

% ==========================================================
\section{Inputs: features and structural encodings}
% ==========================================================
\subsection{Feature sets}
We define three input configurations tailored to synthetic graphs:
\begin{enumerate}[leftmargin=2em]
\item \textbf{Option 0 (structural-only):} degree, $\log(1+\text{degree})$, and LapPE ($k$ eigenvectors).
\item \textbf{Option 1 (random features):} $X\\sim\\mathcal{N}(0,I)$ or one-hot IDs (noting that one-hot can be risky inductively).
\item \textbf{Option 2 (no node features):} structure-only learning via message passing (often very relevant for BC).
\end{enumerate}
For LapPE, we use the normalized Laplacian:
\[
L=I-D^{-1/2}AD^{-1/2},\quad LU=U\\Lambda,\quad X'=[X\\;\\Vert\\;U_k],
\]
and apply random sign flips during training to handle eigenvector sign ambiguity.

% ==========================================================
\section{Models to compare}
% ==========================================================
\subsection{Baselines}
\paragraph{Non-GNN baselines.}
Constant predictor, MLP on $X$, and MLP on $X+$structure.

\paragraph{Shallow GNN baselines.}
GCN / GraphSAGE / GAT / GIN with $L\in\{2,3\}$ (common practice).

\paragraph{Deep GNN baselines (strong references).}
Representative depth-stabilized architectures (e.g., GCNII, APPNP/GPRGNN, JKNet), used to ensure comparisons are meaningful when $L$ is large.
Our focus is on depth scaling and Top-$K$ recovery, which are critical for BC.

\subsection{Proposed: multi-stream routing in GNNs}
\paragraph{Variants.}
We explicitly compare three routing variants: \textbf{HC-GNN (unconstrained)}, \textbf{mHC-GNN (Sinkhorn)}, and \textbf{mHC-lite-GNN (convex combination of permutations)}.

\subsubsection{Multi-stream state}
Each node $i$ carries $n$ parallel streams at layer $l$:
\[
x_i^{(l)}\in\mathbb{R}^{n\times d}.
\]

\subsubsection{Connection maps (dimensions)}
\[
H^{\mathrm{pre}}_{l,i}\in\mathbb{R}^{1\times n},\qquad
H^{\mathrm{post}}_{l,i}\in\mathbb{R}^{1\times n},\qquad
H^{\mathrm{res}}_{l,i}\in\mathbb{R}^{n\times n}.
\]

\subsubsection{Update rule (interpretable form)}
\[
\bar x_i^{(l)}=H^{\mathrm{pre}}_{l,i}x_i^{(l)}\in\mathbb{R}^{1\times d},
\]
\[
m_i^{(l)}=F_{\mathrm{GNN}}\!\left(\bar x_i^{(l)},\{x_j^{(l)}:j\in\mathcal{N}_i\}\right)\in\mathbb{R}^{1\times d},
\]
\[
x_i^{(l+1)}=H^{\mathrm{res}}_{l,i}x_i^{(l)}+(H^{\mathrm{post}}_{l,i})^\top m_i^{(l)}.
\]

\paragraph{How $H^{\mathrm{pre}}$, $H^{\mathrm{post}}$, $H^{\mathrm{res}}$ are built.}
\textbf{Static HC:} $H$ are learned per-layer global parameters shared across nodes.
\textbf{Dynamic HC (optional):} $H$ depend on $x_i^{(l)}$ via a small MLP + nonlinearity (sigmoid/softmax), often after normalization; we start with static-only for robustness and optionally add dynamic components later (static + dynamic is closest to the paper’s recommended setting).

\paragraph{Compute intuition.}
The expensive backbone $F_{\mathrm{GNN}}$ operates on $\bar x_i^{(l)}\in\mathbb{R}^d$ (not $\mathbb{R}^{nd}$).
Overhead mostly comes from multi-stream activations and the small $n\times n$ residual mixing.

% ==========================================================
\section{mHC vs mHC-lite: enforcing stability}
% ==========================================================
\subsection{HC (unconstrained multi-stream baseline)}
We include \textbf{HC-GNN} as the direct multi-stream baseline, using the same update rule but \emph{without} enforcing any constraint on $H^{\mathrm{res}}_{l,i}$.
This isolates the effect of the manifold constraint: HC tests whether routing alone helps depth, while mHC/mHC-lite test whether \emph{constraining} the residual mixing is necessary for stable deep training.

\subsection{mHC (Sinkhorn projection)}
mHC constrains $H^{\mathrm{res}}_{l,i}$ to be approximately doubly-stochastic ($\in B_n$) by:
(i) building a positive matrix (via exponentiation) and
(ii) applying $T$ Sinkhorn-Knopp iterations (row/column normalization).

\subsection{mHC-lite (exact by construction)}
Using Birkhoff--von Neumann:
\[
H^{\mathrm{res}}=\sum_{k=1}^{n!}a_kP^{(k)},\qquad a=\mathrm{softmax}(z),
\]
so $a_k\ge 0$ and $\sum_k a_k=1$ guarantee a convex combination of permutation matrices, hence exact doubly-stochasticity.

\subsection{Scaling issue and approximation}
Since $n!$ grows quickly ($4!=24$, $10!=3{,}628{,}800$), we consider approximations for larger $n$:
selecting a subset of permutations (Top-$M$ or stochastic sampling) and re-normalizing weights, while preserving convex-combination form.

% ==========================================================
\section{Experimental protocol}
% ==========================================================
\subsection{Ablation axes}
We vary:
\[
L\in\{2,4,8,16,32\},\quad n\in\{2,4\},\quad \text{inputs}\in\{\text{struct-only},\text{random},\text{none}\}.
\]
Backbone order: start with GCN (fast, stable), then expand to SAGE/GAT/GIN.

\subsection{Inductive stage (graph family generalization)}
Train on many small graphs (e.g., 100 graphs of 100 nodes), test on larger graphs (e.g., multiple graphs of 1000 nodes), across multiple canonical families:
ER, BA, SBM, WS, RGG (at least five).
We report both ID (train families = test families) and OOD (train families $\neq$ test families) settings to test transfer beyond family-specific style.

% ==========================================================
\section{Reproducibility and reporting}
% ==========================================================
We use fixed random seeds, save configs and checkpoints, and report:
Spearman/NDCG/Precision@K, training curves, runtime and (if possible) memory footprint.

% ==========================================================
\section{Expected outcomes}
% ==========================================================
We expect that naive depth degrades performance without explicit routing mechanisms,
that HC already improves depth scaling,
that mHC/mHC-lite become important as $L$ grows (stability/variance/collapse),
and that Top-$K$ retrieval improves on larger graphs.

% ==========================================================
\section{Repository layout (minimal)}
% ==========================================================
\begin{itemize}[leftmargin=2em]
\item \texttt{configs/} : YAML configs (dataset/model/training/metrics)
\item \texttt{src/data/} : loaders, splits, BC label generation, LapPE
\item \texttt{src/models/} : baselines, deep baselines, HC/mHC/mHC-lite wrappers
\item \texttt{src/train.py} : training loop + logging
\item \texttt{src/eval.py} : ranking and Top-$K$ evaluation
\item \texttt{src/utils/} : seeds, checkpoints, logging
\end{itemize}

\end{document}